\documentclass[11pt]{article}
\usepackage{mathunal-base}
\mathunalfuente{serif}

\renewcommand{\examtitulo}{Ecuaciones Diferenciales (1000007) --- Segundo Examen Parcial}
\renewcommand{\examfecha}{Semestre 02-2015, 26 de Octubre de 2015, Supletorio}

\begin{document}

\examheader

\vspace{4pt}
\noindent
\begin{minipage}[t]{0.62\linewidth}
\textbf{Duración: 1 hora y 40 minutos}

\vspace{8pt}
Nombre y carnet: \dotfill

\vspace{6pt}
Profesor y grupo: \dotfill
\end{minipage}%
\hfill
\begin{minipage}[t]{0.32\linewidth}
\begin{tabular}{|l|c|}
\hline
\multicolumn{2}{|c|}{\textbf{Calificación}} \\ \hline
1. & \\ \hline
2. & \\ \hline
3. & \\ \hline
4. & \\ \hline
\textbf{Total} & \\ \hline
\end{tabular}
\end{minipage}

\vspace{6pt}
\noindent\textbf{Instrucciones.} No está permitido el uso de calculadora ni de cualquier otro tipo de aparato electrónico, incluyendo teléfonos móviles, los cuales deben permanecer apagados durante el tiempo de duración del examen. Tampoco está permitido hacer preguntas a las personas encargadas de la vigilancia.

\vspace{10pt}
\begin{enumerate}
\item{} En los numerales i, ii, iii, iv y v complete según corresponda.
\begin{enumerate}
\item{} (5\%) La posición de un determinado sistema de masa-resorte no amortiguado satisface la ecuación diferencial
\[
\frac{3}{2}u''(t)+k\,u(t) = \operatorname{sen}(2t),
\]
\noindent donde $k$ es una constante positiva. El sistema está en resonancia si el valor de $k$ es: \dotfill.

\item{} (5\%) Si para cualquier solución de la ecuación diferencial con coeficientes constantes
\[
\alpha y''(t)+(2\alpha-1)y'(t)+(\alpha-1)y(t)=0
\]
\noindent se tiene que $\displaystyle\lim_{t\to+\infty}y(t)=0$, entonces la constante $\alpha$ debe estar en el intervalo: \dotfill

\item{} (5\%) Una forma apropiada de una solución particular de la ecuación diferencial
\[
y'''(t)-y'(t) = 3t^2+te^{-t}+2t\cos(t)
\]
\noindent es: \dotfill. (No evalúe constantes)

\item{} (5\%) El mayor intervalo abierto en el que se tiene certeza de que el problema de valor inicial
\[
x(x-2)y''(x)+3xy'(x)+2y(x)=5\ln|7-x|, \qquad y(4)=1,\ y'(4)=2,
\]
\noindent tiene solución única es: \dotfill

\item{} (5\%) Una ecuación diferencial de Cauchy-Euler $ax^2y''(x)+bxy'(x)+cy(x)=0$, donde $a,b,c$ son constantes reales, tiene a $r_1=-1+2i$ como raíz de la ecuación auxiliar. La ecuación diferencial es \dotfill.
\end{enumerate}

\item{}
\begin{enumerate}
\item{} (15\%) Si el Wronskiano $W$ de $f$ y $g$ es $3e^{4t}$ y si $f(t)=e^{2t}$, encuentre $g(t)$.

\item{} (10\%) Si $y_1$ y $y_2$ son soluciones linealmente independientes de la ecuación diferencial $ty''+2y'+te^ty=0$ y si el Wronskiano de $y_1$ y $y_2$ evaluado en 1 es $W(y_1,y_2)(1)=2$, halle $W(y_1,y_2)(5)$.
\end{enumerate}

\item{} Considere la ecuación diferencial lineal
\begin{enumerate}
\item{} (10\%)
\[
a_2x^2\frac{d^2y(x)}{dx^2}+a_1x\frac{dy(x)}{dx}+a_0y(x) = f(x)
\]
\noindent donde $a_0,a_1,a_2$ son constantes reales, con $a_2\neq 0$, y $f$ es una función continua. Pruebe que la sustitución $y(x)=x^m$ transforma la ecuación diferencial dada en una ecuación diferencial con coeficientes constantes. ¿Cuál es la nueva ecuación diferencial?

\item{} (20\%) Halle la solución general de la ecuación diferencial
\[
x^2y''+xy'+y = \sec(\ln x), \qquad x>0.
\]
\end{enumerate}

\item{} Una masa que pesa 8 lbf alarga 6 pulgadas un resorte. Se supone que no hay amortiguamiento y que sobre el sistema actúa una fuerza externa de $8\operatorname{sen}(8t)$ lbf. Si la masa se lleva hacia abajo 3 pulgadas a partir de su posición de equilibrio y luego se libera.
\begin{enumerate}
\item{} (20\%) Determine la posición $x(t)$ de la masa en cualquier instante $t$.
\item{} (6\%) Bosqueje la gráfica de la solución transitoria.
\item{} (4\%) Bosqueje la gráfica de la solución estacionaria.
\end{enumerate}
\noindent (Recuerde que $g$ es 32 pies sobre segundo al cuadrado)
\end{enumerate}

\examfooter
\end{document}
